\begin{table}[t]
\centering
\caption{Cross-validation between declared tool semantics and runtime network evidence.}
\label{tab:semantic-cross-validation}
\begin{tabular}{lrp{0.45\linewidth}}
\toprule
Relation & Repos & Interpretation \\
\midrule
Declared egress risk and observed egress & 5 & Confirmed by runtime evidence \\
Declared egress risk only & 17 & Declared capability was not exercised by this run \\
Observed egress only & 2 & Potential metadata understatement or classifier miss \\
Neither declared nor observed & 11 & No egress evidence in this run \\
\bottomrule
\end{tabular}
\end{table}

\paragraph{Anonymized real-world cases.}
We selected three dynamically executed repositories from the 100-repository
corpus to illustrate the evidence behind Table~\ref{tab:semantic-cross-validation}.
Repo A exposed a code-review tool whose metadata implied GitHub/code repository
and network capability; the invoked tool returned an error containing
\texttt{https://api.github.com/graphql}, while the monitor recorded two denied
connections to \texttt{api.github.com:443}. Repo B exposed a web-search tool
classified as network/browser/database capable; even with invalid generated
arguments, execution attempted an outbound request to a third-party capture API,
which the monitor denied. Repo C exposed a configuration/desktop-control tool
whose called-tool metadata did not classify as network-capable, but the runtime
trace still showed a denied connection to an external browser-support host and
the tool result reflected local configuration and onboarding prompt text. These
cases are anonymized because the purpose is to validate attack-surface evidence,
not to report repository-specific vulnerabilities.
